% \documentclass[margin,line]{C:/Users/LancePC/Obsidian/Main Vault/Job_Market/CV/res}
\documentclass[margin,line]{res}
\usepackage{hyperref}
\usepackage{xcolor}
\usepackage{libertine}
\usepackage{tabularx}


\definecolor{red}{RGB}{192, 8, 52}
\hypersetup{
    colorlinks=true,
    linkcolor=red,
    urlcolor=red,
    citecolor=red
}

\oddsidemargin -.5in
\evensidemargin -.5in
\textwidth=6.0in
\itemsep=0in
\parsep=0in
% if using pdflatex:
%\setlength{\pdfpagewidth}{\paperwidth}
%\setlength{\pdfpageheight}{\paperheight} 

\newenvironment{list1}{
  \begin{list}{\ding{113}}{%
      \setlength{\itemsep}{0in}
      \setlength{\parsep}{0in} \setlength{\parskip}{0in}
      \setlength{\topsep}{0in} \setlength{\partopsep}{0in} 
      \setlength{\leftmargin}{0.17in}}}{\end{list}}
\newenvironment{list2}{
  \begin{list}{$\bullet$}{%
      \setlength{\itemsep}{0in}
      \setlength{\parsep}{0in} \setlength{\parskip}{0in}
      \setlength{\topsep}{0in} \setlength{\partopsep}{0in} 
      \setlength{\leftmargin}{0.2in}}}{\end{list}}


\begin{document}

\name{\moveleft\hoffset\hbox to \textwidth{%
  Lance Gui\hfill
  {\normalfont\small
   \begin{tabular}[b]{@{}r@{\hspace{2em}}r@{}}
     \href{mailto:lance.gui@uni-mannheim.de}{lance.gui@uni-mannheim.de} & \textit{Cell:} 801-897-7578 \\
     \href{https://lancegui.github.io}{\textit{lancegui.github.io}} & \textit{Last Update:} Jun. 2026
   \end{tabular}}%
}}
\begin{resume}

\section{\sc Employment}
{\bf University of Mannheim}, Germany \\
\begin{list1}
\item[] Postdoctoral Researcher, 2026 - 
\end{list1}
\section{\sc Education}
{\bf University of Arizona}, Tucson, Arizona USA\\
%{\em Department of Statistics} 
\vspace*{-.1in}
\begin{list1}
\item[] Ph.D. in Economics, 2020-2026 
\vspace*{.05in}
\item[] M.A. in Economics, 2021
\end{list1}

{\bf Pepperdine University}, Malibu, California USA\\
%{\em Department of Mathematics and Statistics} 
\vspace*{-.1in}
\begin{list1}
\item[] B.S. Mathematics and B.A. Economics, Minor in Computer Science, \textit{Magna Cum Laude}, 2020
\end{list1}

\section{\sc Fields} Industrial Organization, Health Economics

\section{\sc Working Papers}

\textbf{The Last Mile to First Treatment: \\
Search Under Medication Uncertainty} \hfill \href{https://papers.ssrn.com/sol3/papers.cfm?abstract_id=5458595}{\textit{SSRN Link}}\\
\textcolor{red}{\textit{Job Market Paper}} \\

\vspace{-0.7cm}

\noindent {Consumers search for goods whose availability they cannot verify in advance. When travel costs compound this uncertainty, purchases may fail despite demand. I study these frictions in the market for buprenorphine, the leading medication for opioid use disorder. In Washington administrative claims, fewer than half of first-time patients fill their prescriptions. Normalizing each patient’s default pharmacy to zero search cost identifies search costs and preferences: a purchase elsewhere reveals further search was worthwhile. Search frictions explain 70 percent of treatment failures. Guiding patients to available pharmacies captures most of the gains from requiring every pharmacy to carry the drug.}


\textbf{Closing the Opioid Use Disorder Treatment Gap: \\ Expanding Nurse Practitioners’ Prescriptive Authority} \hfill \href{https://ssrn.com/abstract=4774188}{\textit{SSRN link}} \\
\textit{Reject and Resubmit at Journal of Public Economics} \\
\vspace{-0.7cm}

\noindent {Buprenorphine is an effective treatment for opioid use disorder (OUD), yet access remains limited. I study whether expanding buprenorphine prescribing capacity increases treatment access and lowers opioid-related mortality. I exploit the 2016 Comprehensive Addiction and Recovery Act (CARA), which authorized nurse practitioners (NPs) to prescribe buprenorphine for the first time. I combine this federal reform with pre-existing state differences in NP prescribing authority, which determine how strongly CARA expands local access. Granting NPs independent prescribing authority increased active buprenorphine prescribers by 47\%, increased buprenorphine dispensation by 26\%, and reduced opioid-related mortality by 22\%. These effects reflect new provider entry, with no evidence of displacement of existing specialists or opioid treatment programs. Effects are concentrated in previously underserved counties; counties with established access show higher dispensation but no significant mortality decline. I also find suggestive evidence of increased diversion into secondary markets.}

\textbf{Anatomy of Opioid Diversion: Examining Supply-Side Curtailment} \hfill \href{https://ssrn.com/abstract=5044557}{\textit{SSRN link}} \\
With Mo Xiao, Chuan Qin \\ 
\vspace{-0.7cm}

\noindent {Pharmacies should act as gatekeepers of opioid distribution. However, rogue pharmacies divert opioids to non-medical users, worsening the opioid epidemic. We examine the spatial redistribution of opioids dispensed following targeted shutdowns by the Drug Enforcement Administration, employing comprehensive pharmacy-level opioid shipments and hospital diagnoses data in ten U.S. states. The displacement of local opioid shipments after removing a pharmacy unveils the extent of non-medical use: medical users can readily switch to competing local pharmacies, while non-medical users cannot. We develop and estimate a structural model to study medical and non-medical consumers’ substitution patterns facing changes in their consideration sets. We find that 8\% of pharmacies sell non-medical use opioids with a probability greater than 90\% and that over half of pharmacy-dispensed opioids are diverted to non-medical use. Aggressive pharmacy crackdowns, however, drive a substantial portion of non-medical users who lose access to the black market, which may lead to the rise of more dangerous narcotics such as heroin and fentanyl.}


\textbf{Anatomy of a Scandal: \#Batterygate and Consumer Choice} \\
\vspace{-0.7cm}

\noindent  {I link the Twitter activity of over 35,000 individuals to their offline phone purchasing decisions to examine the consumer response (across both existing and potential customers) to the Batterygate scandal. This incident refers to Apple’s 2017 decision to slow down iPhone processors without disclosing the reason, which was to preserve battery life. To address concerns about the parallel trends assumption, I use a novel trend-extrapolated Differences-in-Differences technique that assumes pre-existing trends for the treatment and control groups would have continued without intervention. My findings reveal that Batterygate negatively affected potential consumers but had no impact on the current user base. This suggests that consumer heterogeneity is a crucial factor in determining how customers respond to even uniformly negative media coverage.}

\section{\sc Presentations}
\begin{tabular}{@{}p{5.6in}p{1in}}
     American University, Villanova University, UMass Boston, International Industrial Organization Conference (IIOC), The European Association for Research in Industrial Economics (EARIE) Annual Conference  & 2026 \\
     Association for Public Policy Analysis \& Management (APPAM), Tsinghua University & 2025 \\
     IIOC & 2024 \\
     SUFE‑Jinan Empirical IO Conference, Chinese Economist Society (CES) Annual Conference& 2023 \\
     EARIE Annual Conference, EARIE Inaugural Summer School, Chinese Economist Society (CES) Annual Conference& 2022
\end{tabular}

\section{\sc Teaching}
    University of Arizona, \textit{Sole Instructor} \\
    ECON300 Microeconomics Analysis For Business Decisions Online \vspace*{.05in} \\
    University of Arizona, \textit{Teaching Assistant }\\ 
\hspace*{-8px}
\begin{tabular}{ll}
ECON460 Industrial Organization & ECON462 Firms, Markets, and Competition \\
ECON551 (MBA) Business Strategy & ECON407 Economics of Strategy  \\
ECON508 (Ph.D.) Applied Economic Analysis  & ECON453 Data Analytics and Modeling \\ 
\end{tabular}

\section{\sc Fellowship and Awards} 
Moshe Dror Research Excellence Award \hfill 2025 \\ 
Edward E. Zajac Prize for the Best Third-Year Paper (Honorable Mention) \hfill 2023 \\
Steve Manos Prize for the Best Second-Year Paper (Honorable Mention) \hfill 2022 \\ 
University of Arizona Graduate Fellowship \hfill 2020


\section{\sc Research Grants} 

\begin{tabular}{@{}p{5.6in}p{1in}}
        Lundgren Retail Collaborative Grant (\$2000) with Professor Mo Xiao & 2024 \\
        Center for Management Innovations in Healthcare (CMIH) Research Grant (\$5000) with Professor Mo Xiao &  \\
        CMIH Data Grant (\$2000) with Professor Mo Xiao & \\ 
        Eller Small Grant (\$3000) with Professor Mo Xiao & 2023 \\
\end{tabular}

\section{\sc Research Assistant} Research Assistant to Mo Xiao, University of Arizona \hfill 2023 - 2024 \\
Research Assistant to Eric Hamilton, Pepperdine University \hfill 2016 - 2020

\section{\sc Referee Service} International Journal of Industrial Organization, American Journal of Health Economics

\section{\sc Other Publications} Lee, S.B., \textbf{Gui, X.}, Manquen, M. and Hamilton, E.R., 2019, October. "Use of Training, Validation, and Test Sets for Developing Automated Classifiers in Quantitative Ethnography." In International Conference on Quantitative Ethnography (pp. 117-127).

\noindent Lee, S., \textbf{Gui, X.}, and Hamilton, E. 2020. "Application of AutoML in the Automated Coding of Educational Discourse Data." In Gresalfi, M. and Horn, I. S. (Eds.), The Interdisciplinarity of the Learning Sciences, 14th International Conference of the Learning Sciences (ICLS) 2020, Volume 5 (pp. 2597-2600).

\section{\sc Non-Human Language Skills} R, Julia, Python, MATLAB, C++ {\small (\textit{in order of usage})}

\section{\sc References}

\begin{tabular}{@{}p{1.5in}p{2.2in}p{2in}}
    Professor Mo Xiao & Professor Matthijs Wildenbeest & Professor Hidehiko Ichimura \\
    \href{mailto:mxiao@arizona.edu}{mxiao@arizona.edu} & \href{mailto:wildenbeest@arizona.edu}{wildenbeest@arizona.edu} & \href{mailto:ichimura@arizona.edu}{ichimura@arizona.edu} \\
\end{tabular}




\end{resume}
\end{document}



